\documentclass[11pt]{article}
\usepackage[T1]{fontenc}
\usepackage[ngerman]{babel}
\usepackage{amsmath,amssymb,amsthm}
\usepackage{hyperref}

\title{Eine strukturelle Perspektive auf große Abweichungen in der ABC-Vermutung}
\author{Thomas Hoffbauer}
\date{}

\newtheorem{theorem}{Theorem}
\newtheorem{lemma}{Lemma}
\newtheorem{remark}{Bemerkung}

\begin{document}
	\maketitle
	
	\begin{abstract}
		Wir untersuchen die ABC-Vermutung über eine Analyse der Größe
		\[
		\Delta(a,b,c) := \log c - \log \mathrm{rad}(abc).
		\]
		Numerische Experimente legen nahe, dass große Werte von $\Delta$
		nur bei stark eingeschränkter Primfaktorstruktur auftreten.
		Wir formulieren diese Beobachtung als strukturelles Lemma und diskutieren,
		wie Resultate über S-Unit-Gleichungen und lineare Formen in Logarithmen
		(Baker-Theorie) diese Hypothese stützen können.
	\end{abstract}
	
	\section{Einleitung}
	
	Die ABC-Vermutung besagt, dass für teilerfremde $a,b \in \mathbb{N}$ mit $a+b=c$
	und jedes $\varepsilon>0$ gilt:
	\[
	c \le K_\varepsilon \cdot \mathrm{rad}(abc)^{1+\varepsilon}.
	\]
	
	Eine äquivalente Formulierung betrachtet die Größe
	\[
	\Delta(a,b,c) := \log c - \log \mathrm{rad}(abc).
	\]
	
	Die Vermutung impliziert, dass $\Delta$ nicht wesentlich positiv wachsen kann.
	
	\section{Empirische Beobachtung}
	
	Numerische Experimente (bis in den Bereich mehrerer Millionen) zeigen:
	\begin{itemize}
		\item $\Delta$ ist typischerweise negativ.
		\item Große Werte von $\Delta$ sind selten.
		\item Große $\Delta$ treten nur bei kleiner Anzahl verschiedener Primfaktoren auf.
	\end{itemize}
	
	Diese Beobachtung motiviert das folgende Lemma.
	
	\section{Strukturelles Lemma (Hypothese)}
	
	\begin{lemma}[Strukturelle Enge]
		Für jedes $\delta>0$ existiert eine Konstante $K(\delta)$ mit
		\[
		\Delta(a,b,c) \ge \delta \;\Rightarrow\; \omega(abc) \le K(\delta).
		\]
	\end{lemma}
	
	Wir betonen, dass dieses Lemma derzeit nicht bewiesen ist und als
	leitende Hypothese zu verstehen ist.
	
	\section{Verbindung zu S-Unit-Gleichungen}
	
	Sei
	\[
	S := \{p \text{ prim} : p \mid abc\}, \qquad s := |S|.
	\]
	
	Setze
	\[
	x = \frac{a}{c}, \qquad y = \frac{b}{c}.
	\]
	Dann gilt
	\[
	x + y = 1,
	\]
	wobei $x,y$ S-Einheiten sind.
	
	Nach klassischen Resultaten über S-Unit-Gleichungen (vgl. \cite{EvertseGyory})
	ist die Anzahl solcher Lösungen für festes $S$ endlich.
	
	\section{Baker-Theorie und Größenabschätzungen}
	
	Resultate über lineare Formen in Logarithmen (Baker, Wüstholz) liefern
	explizite Schranken für Lösungen solcher Gleichungen (vgl. \cite{BakerWustholz}).
	
	Insbesondere existieren Konstanten $c_1,c_2>0$, so dass
	\[
	\log c \le c_2 \exp(c_1 s)\, P\, (\log P)^s,
	\]
	wobei $P$ der größte Primteiler von $abc$ ist.
	
	\section{Heuristische Konsequenz}
	
	Unter Verwendung des Primzahlsatzes ergibt sich heuristisch
	\[
	\log \mathrm{rad}(abc) \gtrsim s \log s.
	\]
	
	Damit folgt, dass große Werte von $\Delta$ nur auftreten können,
	wenn $s$ beschränkt ist.
	
	\section{Quantitative Abschätzung (heuristisch)}
	
	Für hinreichend kleine $\delta>0$ führt eine grobe Inversion der Abschätzungen heuristisch auf
	\[
	K(\delta) \lesssim \frac{\log(1/\delta)}{\log\log(1/\delta)}.
	\]
	
	\begin{remark}
		Die vorstehende Abschätzung ist nur als heuristische Größenordnung
		für kleine $\delta$ zu verstehen. Für allgemeines $\delta>0$ ist
		diese Formel in der gegebenen Form nicht sinnvoll, da die Ausdrücke
		$\log(1/\delta)$ und $\log\log(1/\delta)$ nur in einem geeigneten
		Bereich definiert bzw. positiv sind.
	\end{remark}
	
	\section{Diskussion}
	
	Der vorgeschlagene Ansatz interpretiert die ABC-Vermutung als
	Strukturaussage über Primfaktorverteilungen:
	
	Große Abweichungen $\Delta$ scheinen nur bei stark eingeschränkter
	Primstruktur möglich zu sein.
	
	Die Verbindung zu S-Unit-Gleichungen und Baker-Theorie legt nahe,
	dass diese Einschränkung tief in bekannten diophantischen Strukturen
	verankert ist.
	
	\section{Ausblick}
	
	Ein möglicher Zugang zur ABC-Vermutung besteht darin,
	das strukturelle Lemma rigoros zu beweisen oder durch bekannte
	diophantische Methoden zu ersetzen.
	
	\begin{thebibliography}{9}
		
		\bibitem{BakerWustholz}
		A.~Baker und G.~Wüstholz,
		\textit{Logarithmic Forms and Diophantine Geometry},
		Cambridge University Press, 2007.
		
		\bibitem{EvertseGyory}
		J.-H.~Evertse und K.~Győry,
		\textit{Unit Equations in Diophantine Number Theory},
		Cambridge University Press, 2015.
		
	\end{thebibliography}
	
\end{document}